\scanpage{54}
\mspage{3}

On doit bien sûr pour la suite supposer que $W$ satisfait aux conditions équivalentes de la proposition. On prendra p. ex la forme

% slanted
\margin{$W(3)$ Axiome de l'élément final}

\textcircled{3} Si $X$ a un objet final, alors $X$ est $W$-asphérique.

\uline{Corollaire 1} Pour que $f:X\unsure{\text{---}|\longrightarrow}Y$ soit $W$-asphérique, il f. et s. que les $X_{/y}$ soient $W$-asphériques $(y\in\mathrm{Ob}\,Y)$.

Car pour que
\diagram{\begin{tikzcd}[ampersand replacement=\&,column sep=small,row sep=tiny]
X_{/y} \arrow[rr] \arrow[dr,dash] \& \& Y_{/y} \arrow[dl,dash]\\
\& e
\end{tikzcd}}
soit dans $W$, comme $Y_{/y}$ est $W$-asph \inserted{(objet final)}, il f. et s. que $X_{/y}$ le soit aussi.

\deleted{\uline{\unsure{Proposition}}} \uline{Corollaire 2} Si $X$ dans $\mathrm{Cat}$ a un objet initial, il est $W$-asphérique. (En effet, $X\longrightarrow e$ est un homotopisme, et on applique Prop 1)

% slanted
\margin{C'est donc l'énoncé dual de l'axiome \textcircled{3} de l'objet final.}

Je mettrai pour mémoire

\uline{Corollaire 3} Tout homotopisme dans $\mathrm{Cat}$ est dans $W$. Si $f,g:X\rightrightarrows Y$ \deleted{$\mathrm{Cat}$} dans $\mathrm{Cat}$ sont homotopes, alors elles sont égales dans $\mathrm{Hot}_W$, et $f$ est dans $W$ ssi $g$ l'est.

\uline{Proposition 2} : Soit $f\unsure{\in|:}\deleted{\unsure{W}}X\longrightarrow Y$ dans $\mathrm{Cat}$. Si $f\in W$, $f^\circ\in W^\circ$.

On utilise la dém. de I p. 75, à l'aide de la catégorie \unsure{nommée} $\unsure{\Phi|\mathscr{F}}(X)$ $(\mathrm{Ob}\,\Phi(X)=\mathrm{Fl}(X)$, mais avec une notion de \deleted{\unsure{flèches}}\unsure{flèches} dans $\Phi(X)$ un peu \unsure{particulière|particularisée}). On trouve que $f\in W\Longleftrightarrow\Phi(f)\in W$, et de même $f^\circ\in W\Longleftrightarrow\deleted{\unsure{f^\circ\in W}}\Phi(f)\in W$.

NB. Dans la dém. de la prop 2, comme dans celle de la prop. 1 (pour l'implication b) $\Rightarrow$ a)), on n'a utilisé la propriété de localisation \textcircled{2} que sous la forme faible du corollaire, i.e. \deleted{\unsure{comme}} sous la forme $f$ $W$-asphérique $\Rightarrow f\in W$.

% slanted
\margin{donc si on remplace dans cet énoncé $W(2)$ par $W(\unsure{2}\,\mathrm{bis})$.}

\uline{Corollaire} (de la prop. 2) : Soit $W\subset\mathrm{Fl}(\mathrm{Cat})$ satisfaisant l'axiome \textcircled{1} (saturation). Alors la conjonction des axiomes \textcircled{2} \textcircled{3} équivaut à celle des axiomes duaux \textcircled{2$'$} \textcircled{3$'$}.

Il suffit de prouver que $1,2,3\Longrightarrow\unsure{2'|1'},3'$,

% archive sheet number: 54
